\subsection{DECOMPOSITION OF G INTO DOUBLE COSETS}

THEOREM 1. We have G = BWB. The map \(w \mapsto C(w)\) is a bijection from W to the set B\textbackslash G/B of double cosets of G with respect to B.

It is clear that BWB is stable under \(x \mapsto x^{-1}\) , and Lemma 1 shows that it is stable under the product. Since it contains B and N, it is equal to G.

It remains to prove that \(\mathrm{C}(w) \neq \mathrm{C}(w')\) if \(w \neq w'\) . For this, we shall prove by induction on the integer \(q\) the following assertion:

(Aq) If w and \(w'\) are distinct elements of W such that \(l_{S}(w) \geqslant l_{S}(w') = q\) , then \(C(w) \neq C(w')\) .

(For the definition of \(l_{\mathbb{S}}(w)\) , see §1, no. 1.)

This assertion is clear for \(q = 0\) , since then \(w' = 1\) and \(w \neq 1\) , hence \(C(w') = B\) and \(C(w) \neq B\) .

Assume that \(q \geqslant 1\) and that \(w, w'\) satisfy the hypotheses of \((\mathbf{A}_q)\) . There exists \(s \in S\) such that \(sw'\) is of length \(q - 1\) . We have

\[
l _ {S} (w) > l _ {S} \left(s w ^ {\prime}\right)\tag{6}
\]

hence \(w \neq sw'\) . Moreover, \(sw \neq sw'\) ; by formula (3) of §1, no. 1, we have

\[
l _ {\mathrm{S}} (s w) \geqslant l _ {\mathrm{S}} (w) - 1 \geqslant l _ {\mathrm{S}} (s w ^ {\prime}) = q - 1.\tag{7}
\]

By the induction hypothesis, \(C(sw')\) is distinct from \(C(w)\) and from \(C(sw)\) ; from formula (2) it follows that

\[
\mathrm{C} (s w ^ {\prime}) \cap \mathrm{C} (s). \mathrm{C} (w) = \varnothing .\tag{8}
\]

Since \(\mathrm{C}(sw') \subset \mathrm{C}(s).\mathrm{C}(w')\) , we have finally that \(\mathrm{C}(w) \neq \mathrm{C}(w')\) .

Remark. Axiom (T4) was not used in the preceding proof.

